\textbf{Phần 1. Trắc nghiệm lựa chọn}
\setcounterpageref{ex}{0}

\begin{ex} %[1D7V28]
Tính: $14\text{ cm} + 5\text{ cm}$. Kết quả là bao nhiêu xăng-ti-mét?
\choice
{$18\text{ cm}$}
{$20\text{ cm}$}
{$17\text{ cm}$}
{\True $19\text{ cm}$}
\loigiai{
Ta có: $14\text{ cm} + 5\text{ cm} = 19\text{ cm}$.
\textbf{Đáp số:} $19$.\par
\textbf{Đáp án của câu hỏi là D}
}
\end{ex}

\begin{ex} %[1D7V28]
Tính: $18\text{ cm} - 6\text{ cm}$. Kết quả là bao nhiêu xăng-ti-mét?
\choice
{$11\text{ cm}$}
{\True $12\text{ cm}$}
{$13\text{ cm}$}
{$10\text{ cm}$}
\loigiai{
Ta có: $18\text{ cm} - 6\text{ cm} = 12\text{ cm}$.
\textbf{Đáp số:} $12$.\par
\textbf{Đáp án của câu hỏi là B}
}
\end{ex}

\begin{ex} %[1D7V28]
Tính: $20\text{ cm} + 30\text{ cm}$. Kết quả là bao nhiêu xăng-ti-mét?
\choice
{$40\text{ cm}$}
{$60\text{ cm}$}
{\True $50\text{ cm}$}
{$49\text{ cm}$}
\loigiai{
Ta có: $20\text{ cm} + 30\text{ cm} = 50\text{ cm}$.
\textbf{Đáp số:} $50$.\par
\textbf{Đáp án của câu hỏi là C}
}
\end{ex}

\begin{ex} %[1D7V28]
Tính: $70\text{ cm} - 40\text{ cm}$. Kết quả là bao nhiêu xăng-ti-mét?
\choice
{\True $30\text{ cm}$}
{$20\text{ cm}$}
{$40\text{ cm}$}
{$29\text{ cm}$}
\loigiai{
Ta có: $70\text{ cm} - 40\text{ cm} = 30\text{ cm}$.
\textbf{Đáp số:} $30$.\par
\textbf{Đáp án của câu hỏi là A}
}
\end{ex}

\begin{ex} %[1D7V28]
Điền số thích hợp: $12\text{ cm} + \dots\text{ cm} = 18\text{ cm}$.
\choice
{$5\text{ cm}$}
{$7\text{ cm}$}
{$4\text{ cm}$}
{\True $6\text{ cm}$}
\loigiai{
Số cần điền là: $18 - 12 = 6$.
\textbf{Đáp số:} $6$.\par
\textbf{Đáp án của câu hỏi là D}
}
\end{ex}

\begin{ex} %[1D7V28]
Điền số thích hợp: $19\text{ cm} - \dots\text{ cm} = 15\text{ cm}$.
\choice
{$3\text{ cm}$}
{\True $4\text{ cm}$}
{$5\text{ cm}$}
{$2\text{ cm}$}
\loigiai{
Số cần điền là: $19 - 15 = 4$.
\textbf{Đáp số:} $4$.\par
\textbf{Đáp án của câu hỏi là B}
}
\end{ex}

\begin{ex} %[1D7V28]
Đoạn thẳng AB dài 11 cm, đoạn thẳng CD dài 8 cm. Hỏi cả hai đoạn thẳng dài tất cả bao nhiêu xăng-ti-mét?
\choice
{$18\text{ cm}$}
{$20\text{ cm}$}
{\True $19\text{ cm}$}
{$17\text{ cm}$}
\loigiai{
Tổng độ dài hai đoạn thẳng là: $11 + 8 = 19$ (cm).
\textbf{Đáp số:} $19$.\par
\textbf{Đáp án của câu hỏi là C}
}
\end{ex}

\begin{ex} %[1D7V28]
Đoạn thẳng AB dài hơn đoạn thẳng CD ở bài toán trên bao nhiêu xăng-ti-mét?
\choice
{\True $3\text{ cm}$}
{$2\text{ cm}$}
{$4\text{ cm}$}
{$1\text{ cm}$}
\loigiai{
AB dài hơn CD là: $11 - 8 = 3$ (cm).
\textbf{Đáp số:} $3$.\par
\textbf{Đáp án của câu hỏi là A}
}
\end{ex}

\begin{ex} %[1D7V28]
Một sợi dây thừng dài 50 cm. Bác thợ cắt bớt đi 20 cm. Hỏi sợi dây thừng còn lại dài bao nhiêu xăng-ti-mét?
\choice
{$20\text{ cm}$}
{$40\text{ cm}$}
{$29\text{ cm}$}
{\True $30\text{ cm}$}
\loigiai{
Chiều dài sợi dây thừng còn lại là: $50 - 20 = 30$ (cm).
\textbf{Đáp số:} $30$.\par
\textbf{Đáp án của câu hỏi là D}
}
\end{ex}

\begin{ex} %[1D7V28]
Tính giá trị biểu thức: $15\text{ cm} + 4\text{ cm} - 9\text{ cm}$.
\choice
{$0\text{ cm}$}
{\True $10\text{ cm}$}
{$20\text{ cm}$}
{$9\text{ cm}$}
\loigiai{
Thực hiện lần lượt: $15 + 4 = 19$; $19 - 9 = 10$ (cm).
\textbf{Đáp số:} $10$.\par
\textbf{Đáp án của câu hỏi là B}
}
\end{ex}

\begin{ex} %[1D7V28]
Tính giá trị biểu thức: $18\text{ cm} - 8\text{ cm} + 5\text{ cm}$.
\choice
{$14\text{ cm}$}
{$16\text{ cm}$}
{\True $15\text{ cm}$}
{$13\text{ cm}$}
\loigiai{
Thực hiện lần lượt: $18 - 8 = 10$; $10 + 5 = 15$ (cm).
\textbf{Đáp số:} $15$.\par
\textbf{Đáp án của câu hỏi là C}
}
\end{ex}

\begin{ex} %[1D7V28]
Một cuộn ruy băng dài 40 cm, mẹ mua thêm một cuộn ruy băng dài 30 cm. Hỏi mẹ có tất cả bao nhiêu xăng-ti-mét ruy băng?
\choice
{\True $70\text{ cm}$}
{$60\text{ cm}$}
{$80\text{ cm}$}
{$69\text{ cm}$}
\loigiai{
Tổng số ruy băng mẹ có là: $40 + 30 = 70$ (cm).
\textbf{Đáp số:} $70$.\par
\textbf{Đáp án của câu hỏi là A}
}
\end{ex}

\textbf{Phần 2. Trắc nghiệm đúng sai}
\setcounterpageref{ex}{0}

\begin{ex} %[1D7V28]
Xét tính đúng hoặc sai của các khẳng định toán học sau đây:
\choiceTF
{Kết quả của bài toán «Chiếc khăn quàng đỏ dài 80 cm, chiếc nơ trang trí dài 20 cm. Chiếc khăn quàng đỏ dài hơn chiếc nơ bao nhiêu xăng-ti-mét» là $61\text{ cm}$.}
{Một thanh nhôm dài 16 cm. Người ta cắt thành hai đoạn bằng nhau. mỗi đoạn dài $9\text{ cm}$ xăng-ti-mét.}
{\True Đoạn thẳng MN dài 7 cm. Nối tiếp thêm đoạn thẳng NP dài 8 cm và đoạn PQ dài 5 cm (M, N, P, Q thẳng hàng). đoạn MQ dài $20\text{ cm}$ xăng-ti-mét.}
{\True Một cây nến ban đầu cao 20 cm. Sau khi thắp sáng một thời gian, nến còn lại cao 12 cm. cây nến đã bị cháy ngắn đi $8\text{ cm}$ xăng-ti-mét.}
\loigiai{
\begin{itemize}
    \item Ý a) Sai: Kết quả đúng là $60$, không phải $61$. Khăn quàng đỏ dài hơn nơ là: $80 - 20 = 60$ (cm).
\textbf{Đáp số:} $60$.
    \item Ý b) Sai: Kết quả đúng là $8$, không phải $9$. Vì hai đoạn bằng nhau và $8 + 8 = 16$ nên mỗi đoạn dài $8$ cm.
\textbf{Đáp số:} $8$.
    \item Ý c) Đúng: Độ dài đoạn MQ là: $7 + 8 + 5 = 20$ (cm).
\textbf{Đáp số:} $20$.
    \item Ý d) Đúng: Độ dài nến đã cháy là: $20 - 12 = 8$ (cm).
\textbf{Đáp số:} $8$.
\end{itemize}
\begin{itemize}
    \item a) \textbf{Sai}.
    \item b) \textbf{Sai}.
    \item c) \textbf{Đúng}.
    \item d) \textbf{Đúng}.
\end{itemize}
}
\end{ex}

\begin{ex} %[1D7V28]
Xét tính đúng hoặc sai của các bài toán và phép tính sau đây:
\choiceTF
{\True Một con kiến bò từ điểm A đến điểm B dài 9 cm, rồi bò tiếp từ B đến C dài 10 cm. con kiến đã bò được tất cả $19\text{ cm}$ xăng-ti-mét.}
{Kết quả của bài toán «Tìm số tự nhiên $x$ thỏa mãn: $x\text{ cm} + 5\text{ cm} = 17\text{ cm}$» là $13\text{ cm}$.}
{Kết quả của bài toán «Tìm số tự nhiên $y$ thỏa mãn: $20\text{ cm} - y\text{ cm} = 14\text{ cm}$» là $7\text{ cm}$.}
{\True Đoạn thẳng MN dài 7 cm. Nối tiếp thêm đoạn thẳng NP dài 8 cm và đoạn PQ dài 5 cm (M, N, P, Q thẳng hàng). đoạn MQ dài $20\text{ cm}$ xăng-ti-mét.}
\loigiai{
\begin{itemize}
    \item Ý a) Đúng: Tổng quãng đường kiến bò là: $9 + 10 = 19$ (cm).
\textbf{Đáp số:} $19$.
    \item Ý b) Sai: Kết quả đúng là $12$, không phải $13$. Ta có: $x = 17 - 5 = 12$.
\textbf{Đáp số:} $12$.
    \item Ý c) Sai: Kết quả đúng là $6$, không phải $7$. Ta có: $y = 20 - 14 = 6$.
\textbf{Đáp số:} $6$.
    \item Ý d) Đúng: Độ dài đoạn MQ là: $7 + 8 + 5 = 20$ (cm).
\textbf{Đáp số:} $20$.
\end{itemize}
\begin{itemize}
    \item a) \textbf{Đúng}.
    \item b) \textbf{Sai}.
    \item c) \textbf{Sai}.
    \item d) \textbf{Đúng}.
\end{itemize}
}
\end{ex}

\textbf{Phần 3. Trả lời ngắn}
\setcounterpageref{ex}{0}

\begin{ex} %[1D7V28]
Một cái hộp hình chữ nhật có chiều dài 15 cm, chiều rộng 10 cm. Chiều dài hơn chiều rộng bao nhiêu xăng-ti-mét?
\shortans{5}
\loigiai{
\textbf{Đáp số:} $5$\par
Chiều dài hơn chiều rộng là: $15 - 10 = 5$ (cm).
\textbf{Đáp số:} $5$.
}
\end{ex}

\begin{ex} %[1D7V28]
Băng giấy màu xanh dài 25 cm, băng giấy màu đỏ dài 20 cm. Cả hai băng giấy dài bao nhiêu xăng-ti-mét?
\shortans{45}
\loigiai{
\textbf{Đáp số:} $45$\par
Tổng độ dài hai băng giấy là: $25 + 20 = 45$ (cm).
\textbf{Đáp số:} $45$.
}
\end{ex}

\begin{ex} %[1D7V28]
Một cây tre dài 90 cm. Chú thợ cắt bớt đi 30 cm để làm cán chổi. Cây tre còn lại dài bao nhiêu xăng-ti-mét?
\shortans{60}
\loigiai{
\textbf{Đáp số:} $60$\par
Cây tre còn lại dài: $90 - 30 = 60$ (cm).
\textbf{Đáp số:} $60$.
}
\end{ex}

\begin{ex} %[1D7V28]
Ba bạn kiến cùng bò trên một cành cây: kiến đen bò 10 cm, kiến vàng bò 12 cm, kiến lửa bò 8 cm. Hỏi kiến vàng bò nhiều hơn kiến lửa bao nhiêu xăng-ti-mét?
\shortans{4}
\loigiai{
\textbf{Đáp số:} $4$\par
Kiến vàng bò nhiều hơn kiến lửa là: $12 - 8 = 4$ (cm).
\textbf{Đáp số:} $4$.
}
\end{ex}

\begin{ex} %[1D7V28]
Tổng quãng đường bò của kiến đen và kiến lửa ở bài toán trên là bao nhiêu xăng-ti-mét?
\shortans{18}
\loigiai{
\textbf{Đáp số:} $18$\par
Tổng quãng đường là: $10 + 8 = 18$ (cm).
\textbf{Đáp số:} $18$.
}
\end{ex}

\begin{ex} %[1D7V28]
Một đoạn dây điện dài 60 cm, người ta dùng hết 40 cm để nối bóng đèn. Đoạn dây điện còn lại dài bao nhiêu xăng-ti-mét?
\shortans{20}
\loigiai{
\textbf{Đáp số:} $20$\par
Đoạn dây điện còn lại dài: $60 - 40 = 20$ (cm).
\textbf{Đáp số:} $20$.
}
\end{ex}
